\section{Conclusion\label{sec:conclusion}}

A challenging area for further work is to relax and test rationality assumptions. If decision makers are only ``boundedly rational,'' rewards inferred from observed behavior will be distorted by our attempt to rationalize behavior that may not be fully rational.

The ``model-free'' approach to inferring rewards (e.g., \citet{AE:2022}; \citet{kang:2025}) avoids specifying transition dynamics $p$, but implicitly assumes rational expectations: substituting realized transitions for explicit integration effectively ties the agent's subjective beliefs to the data-generating process. When agents hold distorted beliefs (see the tennis example in Section~\ref{sec:ddc}), model-free estimators cannot distinguish misspecified rewards from misspecified beliefs. Model-based methods retain the flexibility to test for such distortions, though at the cost of requiring a correctly specified transition model.

We conclude that the limitations confronting DDC and IRL are not computational
but rather inherent to inferring rewards and beliefs
from data on states and actions, as noted in our discussion of the identification problem in
section 2.4. The solution to the identification problem requires us to impose strong
prior restrictions (assumptions) about rewards and
beliefs.

If these assumptions are wrong, counterfactual predictions from these
models may not be accurate even though the models may provide a good fit to the demonstration data
used to estimate them. However practical experience suggests that even if our assumptions
may not be precisely correct and human decision makers may not be fully rational,
the assumptions used to identify DDC and IRL models appear to be sufficiently good approximations to
produce realistic counterfactual simulations and leverage limited human demonstration data to train
artificially intelligent agents to perform well in tasks previously requiring human expertise.
The rapid growth in the DDC and IRL literatures suggests that despite the limitations
and challenges we have noted, these methods have enormous practical value.
We are likely to see continued rapid progress enabling these methods
to be even more widely used to transform economics, finance and  artificial intelligence.





